\chapter{Attributes for Performance and Safety}
\label{ch:c20-attributes}

\begin{quote}
\emph{C++20 adds three attributes that matter to performance- and correctness-critical code: \texttt{[[likely]]}/\texttt{[[unlikely]]} give branch-probability hints for code layout, \texttt{[[no\_unique\_address]]} lets empty members occupy zero bytes, and \texttt{[[nodiscard("reason")]]} attaches an explanatory message to the discard warning.}
\end{quote}

\section{[[likely]] and [[unlikely]]}
\label{sec:c20-likely}

\begin{lstlisting}[language=C++,caption={Branch-probability hints on if/else and switch}]
int process(int x) {
    if (x > 0) [[likely]] { return fast_path(x); }
    else       [[unlikely]] { return slow_path(x); }
}

int classify(int code) {
    switch (code) {
        case 0: [[likely]]   return ok();
        case 1: [[unlikely]] return rare_error();
        default:             return other();
    }
}
\end{lstlisting}

The standardized form of \texttt{\_\_builtin\_expect}.

\section{What the Branch Hints Actually Do}
\label{sec:c20-branch-hints-effect}

They influence code layout and register/spill decisions, not the hardware predictor directly. A compiler typically lays the likely path inline (fall-through), exiles the unlikely path to a cold section (better i-cache density), biases the conditional jump, and prioritizes the hot path for register allocation.

\begin{lstlisting}[language=C++,caption={The mental model --- cold paths get exiled}]
void handle(Request& r) {
    if (r.malformed()) [[unlikely]] {
        log_and_reject(r);    // moved out-of-line (cold section)
        return;
    }
    serve(r);                 // common path stays inline, cache-dense
}
\end{lstlisting}

\section{When to Use Branch Hints --- and When Not To}
\label{sec:c20-branch-hints-when}

\begin{itemize}
\item Use only where certain and the path is hot (null checks, overflow guards, malformed-input rejection).
\item Do not sprinkle speculatively --- a wrong hint pessimizes by exiling the real hot path.
\item Prefer Profile-Guided Optimization when available; it measures real frequencies.
\item Measure --- the effect is layout-level; confirm via benchmark and assembly.
\end{itemize}

\section{[[no\_unique\_address]]}
\label{sec:c20-no-unique-address}

\begin{lstlisting}[language=C++,caption={An empty member costs zero bytes with [[no_unique_address]]}]
#include <cstddef>

struct Empty {};   // size 1 by default (objects must be addressable)

struct WithoutAttr { Empty e; int value; };           // e >= 1 byte + padding
struct WithAttr    { [[no_unique_address]] Empty e; int value; };  // e may be 0 bytes

static_assert(sizeof(WithoutAttr) >= sizeof(int) + alignof(int));
static_assert(sizeof(WithAttr)    == sizeof(int));    // Empty vanished
\end{lstlisting}

Brings the Empty Base Optimization to member subobjects; no effect on non-empty members.

\section{[[no\_unique\_address]] in Practice: EBO Without Inheritance}
\label{sec:c20-ebo}

\begin{lstlisting}[language=C++,caption={A container storing a (usually empty) comparator at zero cost}]
#include <functional>

template<class T, class Compare = std::less<T>>
class SortedBox {
    T* data_;
    std::size_t size_;
    [[no_unique_address]] Compare comp_;   // empty std::less<T> costs nothing
public:
    bool less(const T& a, const T& b) const { return comp_(a, b); }
};
\end{lstlisting}

How the standard library keeps stateless allocators/comparators/deleters free, via composition instead of private inheritance. MSVC: \texttt{[[msvc::no\_unique\_address]]}.

\section{[[nodiscard]] with a Message}
\label{sec:c20-nodiscard-message}

\begin{lstlisting}[language=C++,caption={[[nodiscard]] with an actionable message}]
struct [[nodiscard("handle this error code")]] ErrorCode { int value; };

[[nodiscard("leaking this handle leaks the resource")]]
Handle acquire();

[[nodiscard("comparison has no effect if discarded")]]
bool empty() const;

void caller() {
    acquire();        // warning includes the message
}
\end{lstlisting}

Applies to a function or a type; converts a generic warning into the consequence and fix.

\section{Attribute Placement and Portability}
\label{sec:c20-attribute-portability}

\begin{lstlisting}[language=C++,caption={Portable spelling for the layout attribute}]
#if defined(_MSC_VER)
  #define NO_UNIQUE_ADDRESS [[msvc::no_unique_address]]
#else
  #define NO_UNIQUE_ADDRESS [[no_unique_address]]
#endif

struct Wrapper {
    NO_UNIQUE_ADDRESS std::less<int> comp;   // zero-cost on all major compilers
    int value;
};
\end{lstlisting}

Unknown standard attributes are ignored (safe incremental adoption); \texttt{[[no\_unique\_address]]} is ABI-affecting, the other two are not.

\section{Professional Insights}
\label{sec:c20-attributes-insights}

\textbf{Apply branch hints only to known-cold error paths and proven-hot loops, and verify with a benchmark.} Prefer PGO, which measures rather than guesses.

\textbf{Reach for \texttt{[[no\_unique\_address]]} on every stateless policy member in generic code} --- EBO via clean composition.

\textbf{Guard \texttt{[[no\_unique\_address]]} for MSVC --- it is ABI-affecting} and not a safe drop-in across a stable binary boundary.

\textbf{Always supply a reason string with \texttt{[[nodiscard]]}} --- on error codes, RAII factories, and pure observers.

\textbf{Treat attributes as safe-to-adopt because unknown ones are ignored} --- except the layout effect of \texttt{[[no\_unique\_address]]}.
